\section{Solution methods and verification}\label{sec:methods}
\subsection{Analytical methods}
For linear problems with constant coefficients the classical tools of elasticity and heat conduction carry over, with only the algebra changed by the coupling.
\begin{itemize}
\item \textbf{Plane-wave and normal-mode analysis.} For unbounded or periodic media, \eqref{eq:disp} gives complete frequency-domain information. Real frequency with complex wavenumber describes spatial decay under steady excitation; real wavenumber with complex frequency describes temporal decay of a free disturbance. The two attenuation measures coincide only for weak damping ($Q^{-1}\ll1$), and the relation involves the group velocity: $\Imag\omega\approx-v_g\Imag k$. Comparing an experiment against the wrong measure is a common source of apparent inconsistency.
\item \textbf{Integral transforms.} For half-spaces, layers and cylinders, Laplace transformation in time and Fourier (or Hankel) transformation in space reduce the problem to ordinary differential equations with exponential solutions, as in Sections~\ref{sec:transient} and~\ref{sec:surface}. Carrying the heat-flux law as the function $G(s)$ lets one derivation serve the whole template.
\item \textbf{Potential methods and eigenfunction expansion.} Helmholtz potentials decouple shear from dilatation in the bulk, and eigenfunction expansion in the transform domain is often used for finite domains \citep{Singh2020a}.
\item \textbf{Wavefront expansions.} Discontinuities propagate along characteristics determined by the high-frequency limit; \eqref{eq:fronts} is an example. Jump amplitudes follow transport equations along the front and decay exponentially in LS-type theories.
\end{itemize}

\subsection{Numerical inversion of Laplace transforms}\label{sec:inversion}
Closed forms in the transform domain are typically inverted numerically. The commonly used algorithms are the Stehfest method \citep{Stehfest1970}, Durbin's Fourier-series method \citep{Durbin1974} (and the closely related method of Dubner and Abate \citep{DubnerAbate1968}), the Honig--Hirdes method \citep{HonigHirdes1984}, and contour-deformation methods such as Talbot's \citep{Talbot1979}; Abate and Whitt \citep{AbateWhitt2006} review the common framework. Applications to thermoelastic problems often report numbers from one of these methods without a check. Four implementations were therefore benchmarked on four transforms with known inverses: Talbot, Stehfest, de Hoog et al.\ \citep{deHoog1982} and Cohen acceleration \citep{Cohen2007}, all as implemented in a standard multiprecision library with default parameters at 15 digits. The test transforms represent the features of thermoelastic responses: a smooth decay, an oscillation, a discontinuous step with a delay, and a diffusive response (\Cref{tab:bench}).
\begin{table}[ht]\centering\small
\caption{Maximum absolute error of four numerical Laplace inversion methods on test transforms (evaluated at $t=0.5,2,5,10$; for the step at $t=0.5,0.95,1.05,2,5$). Default parameters, 15 digits. ``$4\times10^{193}$'' indicates a numerical blow-up.}\label{tab:bench}
\input{tables/laplace_bench.tex}
\end{table}
Four points deserve recording. (i) For smooth and diffusive responses all four methods are essentially exact ($<10^{-12}$), which is why published diffusion-dominated results rarely depend on the method. (ii) For the step with a delay, Talbot's method \emph{fails catastrophically} at times before the delay ($t=0.95<1$): its contour, which extends into the left half-plane, is multiplied by $e^{-s}$ and grows without bound. A Cattaneo front is exactly this situation, since its transform contains $e^{-x/c_T\cdot s}$. \cref{sec:transient} handles this by setting the solution to zero ahead of the known front and inverting only for $t>x/c_T$. (iii) Stehfest and Cohen err by order $0.2$ near the discontinuity, smearing the front and ringing, whereas de Hoog gives $1.4\times10^{-6}$ at the points examined. (iv) For oscillatory functions the default Talbot setting is inadequate at $t=10$, with error $0.04$, while de Hoog and Cohen stay accurate. Because all methods can be tuned, through the number of terms or the contour parameter, the table is not a ranking; it shows that the choice of method and the checks matter. The recommendation is to invert with two independent methods, to compare against a known limit, and to treat results near discontinuities as unreliable beyond the known front position.

\subsection{Time-domain numerics}
For problems that do not admit transform solutions (heterogeneous, nonlinear, complex geometry) finite-element and finite-difference methods in space and time are used. Two issues belong to hyperbolic thermal models. Schemes built for parabolic heat equations, implicit and numerically dissipative, smear the thermal front and so cannot represent the finite-speed effect that motivates the model. Explicit or high-resolution schemes with controlled dissipation are the appropriate choice, as are space--time discontinuous Galerkin formulations. The second issue is that the time step must resolve the relaxation time ($\Delta t\ll\tau_q$) and meet a Courant--Friedrichs--Lewy (CFL) condition set by the \emph{largest} wave speed, $\max(c_1,c_T)$ and not $c_1$. When $c_T\gg c_1$, as in many metals, the thermal front governs the stability bound. Memory-dependent and fractional operators need stored history and, without optimization, cost $O(N^2)$ in the number of time steps. Short-memory approximations and fast convolution cut this down, and should come with an error estimate. Finite-element implementations with memory-dependent heat transfer in fiber-reinforced cylinders \citep{Sur2019} and rotating nano-beams \citep{Sur2025} illustrate the approach.

\subsection{A verification discipline}
The checks below catch most errors in new thermoelastic wave calculations, and authors should state which of them they performed.
\begin{enumerate}
\item \textbf{Reduce to a known limit}: set the coupling to zero (elastic and thermal problems separately), and set the lags to zero (Fourier).
\item \textbf{Check the dispersion relation symbolically} against the field equations, as in \cref{app:disp}; sign and factor errors in the coupling terms are detected immediately.
\item \textbf{Check energy}: for conservative members of the family (GN-II) the energy is constant; for the others, the energy balance \eqref{eq:energy} must hold with the stated dissipation; for interfaces, the net flux must be continuous (\cref{sec:interfaces}).
\item \textbf{Check high- and low-frequency asymptotics}: adiabatic and isothermal speeds, front speeds (e.g.\ \eqref{eq:fronts}), classical diffusive scaling.
\item \textbf{Check branches}: track roots by continuity and ensure that Taylor truncation has not introduced spurious roots.
\item \textbf{Report dimensionless groups}: $\varepsilon$, $\omega/\omega_*$, $\omega\tau$, $c_T/c_1$, $\xi\omega/c_1$, so that results can be compared across studies.
\end{enumerate}


\begin{summarybox}[Summary]
Solution methods: analytical separation for simple geometries, Laplace/Fourier inversion for transients, numerical root finding for dispersion and secular equations. Verification discipline (combining the checks of this section with those of the interface section): reduce to the elastic limit, check Snell invariants, verify net-flux continuity (not mode-flux sums), test boundary conditions by substitution, state the flux definition for lagged and memory models. Report dimensionless groups $\varepsilon$, $\omega/\omega_*$, $\omega\tau$, $c_T/c_1$, $\xi\omega/c_1$ for comparability.
\end{summarybox}